% META: {"id":"1NSI-PM-RE-C3-CORRIGE","chapitre":"1NSI-PROJET-METHODES","type_objet":"corrige","exercice_ref":"1NSI-PM-RE-C3","capacites":["BO-PREAMBULE-COMPETENCES-METHODE"],"status":"approved","origin":"needs_review","status_history":["needs_review","audited","accepted","approved"],"release_acceptance":"RELEASE_OWNER_BATCH_ACCEPTANCE","acceptance_closure_digest":"sha256:9b3ccf9a81c5520fbb7e03b7b2d3e7bf2057a02834b6d908bf3be5f49f3b553f","acceptance_receipt_digest":"sha256:4fad86f0282e0fa4e0419cca1ad6379ae7af5a280c04a4a90880f90a1c044242"}

\begin{corrige}{1NSI-PM-RE-C3-EX1}
\textbf{1.} Le programme lève une \lstinline{TypeError}, avec le message
\lstinline{can only concatenate str (not "int") to str}.

\textbf{2.} L'expression fautive est \lstinline{"... tu as " + age + " ans."} :
l'opérateur \lstinline{+} entre chaînes de caractères \textbf{concatène} du texte, mais
\lstinline{age} est un entier (\lstinline{int}), pas une chaîne. Python refuse de mélanger
directement un \lstinline{int} et un \lstinline{str} avec \lstinline{+}.

\textbf{3.} Deux corrections possibles :
\begin{python}
def message_bienvenue(nom, age):
    return "Bonjour " + nom + ", tu as " + str(age) + " ans."

def message_bienvenue_fstring(nom, age):
    return f"Bonjour {nom}, tu as {age} ans."
\end{python}
La première convertit explicitement \lstinline{age} en chaîne avec \lstinline{str(...)} ;
la seconde utilise une \emph{f-string}, qui effectue cette conversion automatiquement.
\end{corrige}
